\section{A uniform finite-description realization}\label{sec:uniform52}
The lower bounds allow arbitrary real rows for free. A finite-group
upper realization can instead be priced in a conventional bit model
when the matrices and signal are rational.

\begin{theorem}[Priced rational finite-group realization]
\label{thm:uniform52}
Suppose $\rho$ and the command matrices are rational and $G$ is finite.
There is one horizon-independent exact realization with
$M\le2D|G|$ command-state labels and deterministic command updates.
Given the finite orbit and its permutation tables, the latter occupy
$O(|\calA|M\log M)$ bits. Let $L$ bound the bit length of the
numerators and denominators of the $MD$ terminal Bernoulli
probabilities. Decoder descriptions use $O(MDL)$ bits. During command
processing the persistent register uses $O(\log M)$ bits; exact
terminal sampling can be performed with $O(\log M+L)$ workspace and
fewer than $2L$ expected fair random bits per query. Its worst-case
sampling time is not bounded. These bounds are independent of the
horizon, and the same tables answer at any stopping epoch.
\end{theorem}
\begin{proof}
Enumerate the finite orbit $\{gx:g\in G,x\in X\}$ using rational
arithmetic and duplicate testing. Breadth-first closure under the
finite alphabet terminates because the orbit is finite. Store its
point as the state; every command permutes this set, and the seed
initialization is deterministic. Each label $z$ has coordinate
probabilities $(1+\rho z_j)/2$, which are rational and legal because
$\|z\|=1$. They give exactly the desired numerical interface.

The table counts follow by writing one successor index per command
and state, and one rational per decoder entry. To sample a rational
probability $p/q\in[0,1]$ in lowest terms, handle $p=0,q$ without
randomness; otherwise let $\ell=\lceil\log_2q\rceil\le L$.
Draw $\ell$ fair bits to obtain an integer uniformly in
$\{0,\ldots,2^\ell-1\}$. Reject and repeat when it is at least $q$;
otherwise output success precisely when it is below $p$. Acceptance
probability is greater than $1/2$ (or one for a power of two), so the
expected number of trials is less than two. Counters and comparisons
use $O(L)$ bits of workspace. There is no finite worst-case number
of trials. Construction of the tables is a separate finite computation,
not charged to the persistent label count. No clock is used by the
updates or decoder, so all horizon and stopping-epoch assertions follow.
\end{proof}

The finite-orbit enumeration is a terminating synthesis under the
finite-group hypothesis, not an assertion that group finiteness has
been decided by that enumeration. Heights of generated orbit points
and the construction time are not bounded merely by the atomic width.
The theorem thus distinguishes label memory, description length,
arithmetic, randomness and preprocessing. It does not claim a uniform
bit-cost realization for arbitrary prescribed real transition rows.
